\phantomsection\label{sec:kitti-extended}
The extended observed-first experiments improved all four selected
validation scores beyond their three-epoch counterparts, but none
reached its line-specific reference
(Table~\ref{tab:k-extended-comparison}). The comparison separates
detector adaptation on an unchanged upsampled input from comparison
with the original observation. Thus the frozen observed-first rows
measure the starting point for adaptation, whereas the original-$N$
and sparse-$M$ rows are distinct input controls. The original-$N$ row
also serves as the complete-observation reference for Line~B.

Frozen denotes official detector weights. Extended denotes the
validation-selected best checkpoint from the schedules in
Table~\ref{tab:k-extended-stages}, with the PU-GCN checkpoint unchanged.
Baseline training budgets were not extended.

\begin{table}[!t]\centering\setlength{\tabcolsep}{4pt}
\begin{tabular}{lL{0.43\linewidth}rr}\toprule
Line & Input / detector weights & PointRCNN & CenterPoint\\\midrule
\tableinput{tables/kitti_extended_comparison.tex}
\bottomrule\end{tabular}
\caption{Detector-adaptation comparison: Car Moderate 3D AP$_{R40}$ on
3,769 validation frames. Bold marks extended-schedule results,
not the highest scores across input controls.}
\label{tab:k-extended-comparison}\end{table}

Relative to the same observed-first input with frozen detector weights,
the extended PointRCNN results gained 7.32 AP points in Line~A and
21.52 in Line~B. The corresponding CenterPoint gains were 12.76 and
18.81 points. The additional gains over the earlier three-epoch
adaptation were smaller: 0.58 and 1.83 for PointRCNN, and 1.26 and
1.24 for CenterPoint. Figure~\ref{fig:k-convergence-gains} separates
these recovery gains from the remaining baseline deficits, so a large
improvement from a degraded starting point is not mistaken for an
advantage over the original input.

\thesisfigure{k_convergence_gains}{Recovery and remaining deficits after extended adaptation}
{Car Moderate 3D AP$_{R40}$ differences on 3,769 validation frames.
(a) Extended-best minus frozen observed-first and minus three-epoch
observed-first results. (b) Extended-best minus the original-$N$ frozen
baseline and minus the line-specific reference. The latter is the same
original baseline in Line~A and the three-epoch adapted sparse-$M$
baseline in Line~B; the two markers therefore coincide in Line~A.
Each point comes from one completed schedule and its recorded comparison,
not an average across independent seeds.}{fig:k-convergence-gains}

For Line~A, the remaining original-baseline deficits were 10.36 AP points
for PointRCNN and 3.31 for CenterPoint. In Line~B, the extended results
remained 11.28 and 5.15 points below the adapted sparse controls.
They were also below the \emph{unadapted} sparse controls, by 8.53 and
1.70 points, respectively. Compared with the complete original
observation, the Line~B deficits were larger still, at 24.90 and
16.38 points. The recovery therefore does not depend on choosing only
one unfavourable starting comparator.

The stage summaries in Table~\ref{tab:k-extended-stages} establish the
scope of the plateau assessment. PointRCNN selected RPN epoch~14 in
both lines; the associated RCNN selections were epoch~2 in Line~A and
epoch~19 in the 24-epoch Line~B schedule. The four final PointRCNN
trajectories contain 78 epoch evaluations in total. CenterPoint selected
epoch~12 in each of its two 12-epoch schedules, adding 24 evaluations.
These 102 checkpoint evaluations are repeated measurements of the
completed training trajectories, not 102 independent training runs.
PointRCNN RPN-stage AP belongs to its stage-specific detection
evaluation and is not an additional final RCNN result.

\begin{table}[!t]\centering\setlength{\tabcolsep}{4pt}
\begin{tabular}{lllrrrr}\toprule
Detector & Line & Stage & \shortstack{Total\\epochs} &
\shortstack{Best\\epoch} & \shortstack{Best\\AP} &
\shortstack{Final\\count}\\\midrule
\tableinput{tables/kitti_extended_stages.tex}
\bottomrule\end{tabular}
\caption{Completed schedules and selected checkpoints. Final count denotes
consecutive terminal epochs without improvement exceeding the running
reference by 0.2 AP points (Section~\ref{sec:fc-4-11}).}
\label{tab:k-extended-stages}\end{table}

All six terminal counts in Table~\ref{tab:k-extended-stages} meet the
threshold of three. A numerical maximum at the last epoch does not
contradict this threshold-based plateau rule. CenterPoint Line~A improved from the last
threshold-reference score of 75.7998 at epoch~8 to 75.9661 at epoch~12;
Line~B improved from 62.8977 at epoch~9 to 62.9007 at epoch~12.
These increases of 0.1663 and 0.0030 AP points did not reset the counter.
For PointRCNN Line~B, Figure~\ref{fig:k-convergence-terminal} shows
the six recorded terminal RCNN evaluations. The selected epoch~19
scored 57.0485, whereas epoch~24 scored 55.5902 and all five intervening
or final evaluations remained below the selected value. The last
checkpoint was therefore not the selected checkpoint.

\thesisfigure{k_convergence_terminal}{Terminal validation after the selected PointRCNN checkpoint}
{PointRCNN Line~B Car Moderate 3D AP$_{R40}$ for epochs~19--24
of the completed 24-epoch RCNN schedule, with RPN epoch~14 fixed.
Blue marks the selected RCNN epoch; grey marks subsequent epochs.
The panel shows the complete terminal six-epoch window tabulated in
the experiment record, rather than the full training trajectory.
Every point evaluates all 3,769 validation frames; the connecting
segments join adjacent recorded epochs and do not estimate missing
earlier values.}{fig:k-convergence-terminal}

This extension supplies the validation-progress evidence missing from
the initial three-epoch study. It narrows the explanation for the
remaining deficit: stopping after three epochs alone does not account
for all of the loss under the tested schedules. It does not establish
a global optimum, exclude improvement with another optimiser or input
design, or remove the unequal-budget and validation-selection limits.
The earlier all-class tables retain their original three-epoch
checkpoints; their Pedestrian and Cyclist values are not reassigned
to these Car-selected models.
